\documentclass{article}
\usepackage{amsmath,amssymb,hyperref}
\begin{document}
\title{An optimal five-dimensional second-order anticoherent subspace of 24 qubits}
\author{Catalogue calculation}
\date{30 October 2026}
\maketitle
\newcommand{\D}[1]{|D_{24}^{(#1)}\rangle}
\paragraph{Result.} The maximal dimension is $K(24,2)=5$. The basis below is algebraic, specified by a unique root in a rational box; no short radical expression is asserted.

\paragraph{Reflection prescription.} Let $R\D{r}=\D{24-r}$, with no complex conjugation. Set
\begin{align*}
u_1&=\sum_{\ell=0}^{4}a_\ell\D{5\ell}, & u_5&=Ru_1,\\
u_2&=\sum_{\ell=0}^{4}b_\ell\D{1+5\ell}, &u_4&=Ru_2,\\
u_3&=c_2(\D{2}+\D{22})+c_7(\D{7}+\D{17})+c_{12}\D{12}.
\end{align*}
Here $a_\ell=c_{5\ell}$ and $b_\ell=c_{1+5\ell}$ in terms of the global Dicke amplitudes.
Write each Dicke coefficient as $\sqrt{\binom{24}{r}}\,\xi_{a,r}$.
The independent coefficients are $\xi_{1,5\ell}$ and $\xi_{2,1+5\ell}$ for $0\le\ell\le4$, and $\xi_{3,2},\xi_{3,7},\xi_{3,12}$. All others follow by reflection. The imaginary parts of the first coefficient in each of the three independent vectors are fixed to zero. The remaining 23 real coordinates are denoted by $x$.

\paragraph{Integer quadratic equations.} Define $h_0(r)=1$, $h_1(r)=24-2r$, and $h_2(r)=3(24-2r)^2-624$. For $r^{\underline q}=r(r-1)\cdots(r-q+1)$, put
\[
P_{ab}^{q,p}=\sum_{r=q}^{24}\binom{24}{r}r^{\underline q}h_p(r)
\overline{\xi_{a,r-q}}\xi_{b,r}-\delta_{q0}\delta_{p0}\delta_{ab}.
\]
The real and imaginary parts of $P^{0,p}$ for $p=0,1,2$, $P^{1,p}$ for $p=0,1$, and $P^{2,0}$ are precisely the Gram and rank-one/two moment conditions, up to nonzero scalar factors. After the support and reflection substitutions, there are 42 nonzero real polynomial equations. Their coefficient matrix over the rational numbers has rank 23. The selected 23 equations $F(x)=0$ and the exact rational reconstruction of all 42 equations are listed in \texttt{certificate.json}. Thus no discarded equation is assumed to follow merely from numerical agreement.

\paragraph{Exact coefficient definition.} The coordinates are the unique real root of $F$ in
\[
B=\{x:\|x-x^{(0)}\|_\infty\le10^{-60}\},\qquad
x_i^{(0)}=n_i/10^{80}.
\]
All 23 integer numerators $n_i$ are in the certificate. A rational matrix $A$, also supplied there, satisfies
\begin{align*}
\|I-AF'(x^{(0)})\|_\infty&<3\times10^{-64},\\
\|AF(x^{(0)})\|_\infty&<5\times10^{-81},\\
\|A\|_\infty&<31/100,\\
\|F'(x)-F'(y)\|_\infty&\le7647310848\|x-y\|_\infty.
\end{align*}
Consequently $G(x)=x-AF(x)$ has contraction factor less than $2.4\times10^{-51}$ on $B$, and maps $B$ into a concentric box of radius less than $5\times10^{-81}$. The contraction theorem gives a unique fixed point in $B$. The first inequality implies that $AF'(x^{(0)})$, and hence $A$, is invertible, so the fixed point is a zero of $F$. Every bound is checked using integer and rational arithmetic; the initial numerical search is not the existence proof.

\paragraph{Anticoherence of the entire span.} The moment equations give
\[
\langle u_a|u_b\rangle=\delta_{ab},\qquad
\langle u_a|J_i|u_b\rangle=0,\qquad
\left\langle u_a\left|\frac{J_iJ_j+J_jJ_i}{2}\right|u_b\right\rangle
=52\,\delta_{ij}\delta_{ab}.
\]
Here $j=12$ and $j(j+1)/3=52$. Equivalently,
\[
R_2(|u_a\rangle\langle u_b|)=\delta_{ab}I_3/3\qquad(1\le a,b\le5).
\]
These include all cross terms, so every normalized complex superposition in the span is second-order anticoherent.

\paragraph{Optimality.} The established auxiliary-spin-two inertia bound is $5k\le N+5$ for any second-order AC subspace with $N\ge4$. At $N=24$ this gives $5k\le29$ and therefore $k\le5$. The construction attains this upper bound, proving $K(24,2)=5$ without restricting competing subspaces to the cyclic or reflection ansatz.

\paragraph{Scope.} This is one exact optimal instance of the cyclic-residue prescription. It does not prove the prescription has solutions, or attains the maximum, for every $N$ and $k$.
\end{document}
